\chapter{Boolean Analysis --- Parameters}
%%% generated by the blueprint pipeline from TCSlib.BooleanAnalysis.Hypercontractivity.RandomVariables.Sharp.Parameters
\section{Overview}
Hyperbolic-function formulas describe the sharp radius and its biased two-point coefficient. Monotonicity in atom mass and comparison with the balanced radius give the scalar bounds used for finite laws.

\section{Declarations}

\begin{theorem}[The sharp radius lies between zero and one]
\lean{BooleanAnalysis.Hypercontractivity.SharpDiscrete.radius_pos_lt_one}
\leanok
\uses{BooleanAnalysis.Hypercontractivity.sharpDiscreteRadius}
\difficulty{4}
Let $q>2$ and $0<\lambda<1/2$. Then $0<\rho_*(q,\lambda)<1$. Here $\rho_*(q,\lambda)=\sqrt{\sinh(\ell/q)/\sinh(\ell/(q/(q-1)))}$ with $\ell=\log((1-\lambda)/\lambda)$.
\end{theorem}

\begin{theorem}[Hyperbolic formula for the bias coefficient]
\lean{BooleanAnalysis.Hypercontractivity.SharpDiscrete.bias_coefficient_eq_sinh}
\leanok
\difficulty{5}
Let $p\in\bbr$ and $0<\alpha<\beta$. Then $(\alpha^p\beta^{2-p}-\alpha^{2-p}\beta^p)/(\alpha^2-\beta^2)=\sinh((p-1)\log(\beta/\alpha))/\sinh(\log(\beta/\alpha))$.
\end{theorem}

\begin{theorem}[The bias coefficient lies between zero and one]
\lean{BooleanAnalysis.Hypercontractivity.SharpDiscrete.bias_coefficient_pos_lt_one}
\leanok
\uses{BooleanAnalysis.Hypercontractivity.SharpDiscrete.bias_coefficient_eq_sinh}
\difficulty{4}
Let $1<p<2$ and $0<\alpha<\beta$. Then $0<(\alpha^p\beta^{2-p}-\alpha^{2-p}\beta^p)/(\alpha^2-\beta^2)<1$.
\end{theorem}

\begin{theorem}[The bias coefficient is the squared radius]
\lean{BooleanAnalysis.Hypercontractivity.SharpDiscrete.bias_coefficient_eq_radius_sq}
\leanok
\uses{BooleanAnalysis.Hypercontractivity.SharpDiscrete.bias_coefficient_eq_sinh, BooleanAnalysis.Hypercontractivity.SharpDiscrete.bias_coefficient_pos_lt_one, BooleanAnalysis.Hypercontractivity.sharpDiscreteRadius}
\difficulty{5}
Let $1<p<2$, $0<\alpha<\beta$, and $\alpha^p+\beta^p=1$. Then $(\alpha^p\beta^{2-p}-\alpha^{2-p}\beta^p)/(\alpha^2-\beta^2)=\rho_*(p/(p-1),\alpha^p)^2$. Here $\rho_*(q,\lambda)=\sqrt{\sinh(\ell/q)/\sinh(\ell/(q/(q-1)))}$ with $\ell=\log((1-\lambda)/\lambda)$.
\end{theorem}

\begin{theorem}[Concavity of the hyperbolic tangent]
\lean{BooleanAnalysis.Hypercontractivity.SharpDiscrete.concaveOn_tanh_nonneg}
\leanok
\difficulty{4}
The hyperbolic tangent $x\mapsto\tanh x$ is concave on $[0,\infty)$.
\end{theorem}

\begin{theorem}[Monotonicity of a hyperbolic sine ratio]
\lean{BooleanAnalysis.Hypercontractivity.SharpDiscrete.sinh_ratio_antitoneOn}
\leanok
\uses{BooleanAnalysis.Hypercontractivity.SharpDiscrete.concaveOn_tanh_nonneg}
\difficulty{5}
For $0\le a\le1$, the function $u\mapsto\sinh(au)/\sinh u$ is nonincreasing on $(0,\infty)$.
\end{theorem}

\begin{theorem}[Monotonicity of the radius in atom mass]
\lean{BooleanAnalysis.Hypercontractivity.SharpDiscrete.radius_mono_bias}
\leanok
\uses{BooleanAnalysis.Hypercontractivity.SharpDiscrete.sinh_ratio_antitoneOn, BooleanAnalysis.Hypercontractivity.sharpDiscreteRadius}
\difficulty{5}
Let $q>2$ and $0<\lambda\le\gamma<1/2$. Then $\rho_*(q,\lambda)\le\rho_*(q,\gamma)$. Here $\rho_*(q,\lambda)=\sqrt{\sinh(\ell/q)/\sinh(\ell/(q/(q-1)))}$ with $\ell=\log((1-\lambda)/\lambda)$.
\end{theorem}

\begin{theorem}[Scaling the hyperbolic sine]
\lean{BooleanAnalysis.Hypercontractivity.SharpDiscrete.sinh_scale_le}
\leanok
\difficulty{4}
Let $0\le a\le1$ and $u\ge0$. Then $\sinh(au)\le a\sinh u$.
\end{theorem}

\begin{theorem}[Comparison with the balanced radius]
\lean{BooleanAnalysis.Hypercontractivity.SharpDiscrete.radius_sq_le_balanced}
\leanok
\uses{BooleanAnalysis.Hypercontractivity.SharpDiscrete.radius_pos_lt_one, BooleanAnalysis.Hypercontractivity.SharpDiscrete.sinh_scale_le, BooleanAnalysis.Hypercontractivity.sharpDiscreteRadius}
\difficulty{4}
Let $q>2$ and $0<\lambda<1/2$. Then $\rho_*(q,\lambda)^2\le1/(q-1)$. Here $\rho_*(q,\lambda)=\sqrt{\sinh(\ell/q)/\sinh(\ell/(q/(q-1)))}$ with $\ell=\log((1-\lambda)/\lambda)$.
\end{theorem}
